\section{Local estimation and paired borrowing}
In this section we define the reference target, the augmentation path and the certified borrowing rule. We work conditionally on an independent training sample and treat the verification probability as known; Section~S8 covers estimated probabilities.

\subsection{The reference distribution and augmentation path}
We observe $O=(W,R,RY)$, where $W$ includes $X$ and $S$. Assume $R\perp Y\mid W$ and initially treat $\pi(W)=p e(W)$ as known, with $e$ bounded above and away from zero. Let $Z_y=\ind(Y\leq y)$. For a bounded nonnegative kernel $K_h(X-x_0)=K\{(X-x_0)/h\}$, put
\begin{equation}\label{eq:target}
 \mu_h=\E K_h,\quad v_h=h^d,\quad f_h=\mu_h/v_h,\quad
 F_h(y)=\frac{\E(K_hZ_y)}{\mu_h}.
\end{equation}
Equation~\eqref{eq:target} defines the inferential target at the declared resolution $h$. The point target $F_0(y\mid x_0)$ additionally requires localisation-bias control.

We use an independent training sample that produces bounded candidate distribution functions $q_y(X)$ and $m_y(W)$. The anchor omits the surrogate from its outcome regression. Define $D_y=m_y-q_y$ and $a_{\lambda,y}=q_y+\lambda D_y$, $0\leq\lambda\leq1$. The augmented signal has the linear representation
\begin{equation}\label{eq:signal}
 U_{\lambda,y}=q_y+\frac R\pi(Z_y-q_y)
       +\lambda\left(1-\frac R\pi\right)D_y
       =U_{0,y}+\lambda B_y.
\end{equation}
Equation~\eqref{eq:signal} shows that the correction enters linearly and that, conditional on the training sample, $\E(U_{\lambda,y}\mid W)=\E(Z_y\mid W)$ and $\E(B_y\mid W)=0$. Candidate accuracy changes covariance, not the target.

On a threshold grid $\Yset_J=\{y_1,\ldots,y_J\}$, form
\begin{equation}\label{eq:psi}
 \Psi=\sqrt{p/v_h}\,K_h
       \binom{U_0-F_h}{B},\qquad
 \Omega=\E(\Psi\Psi^\top),\qquad H_\lambda=(I_J,\lambda I_J).
\end{equation}
Equation~\eqref{eq:psi} stacks the anchor score and the correction, and the covariance path is quadratic:
\begin{equation}\label{eq:path}
 \Sigma_\lambda=H_\lambda\Omega H_\lambda^\top
 =\Omega_{00}+\lambda(\Omega_{01}+\Omega_{10})+\lambda^2\Omega_{11}.
\end{equation}
For the evaluation sample $\se$, the kernel ratio admits the exact representation derived in Section~S3, and $\wh f_{h,\se}=\PP_{\se}K_h/v_h$; a zero denominator gives the full $[0,1]$ distribution band.

\subsection{The uncertainty criterion}
For $G_\lambda\sim N(0,\Sigma_\lambda)$, let
\begin{equation}\label{eq:gain}
 c_\Omega(\lambda)=\operatorname{Quantile}_{1-\alpha}\norm{G_\lambda}_\infty,
 \qquad g_\Omega(\lambda)=c_\Omega(0)-c_\Omega(\lambda).
\end{equation}
Equation~\eqref{eq:gain} defines the object of the selector: the reduction in the Gaussian simultaneous critical value at the same resolution, sample size and confidence level. A fixed positive band shape is allowed by rescaling both score blocks before forming $\Omega$; the shape is common to all weights. The trace rule minimises $\tr(\Sigma_\lambda)$ instead.

\begin{proposition}[Separation of the criteria]\label{prop:separation}
There is a valid distributional augmentation path for which the trace minimiser $25/41$ reduces $\tr(\Sigma_\lambda)$ from $0.340$ to $0.318$, but increases its 95\% simultaneous critical value from $0.983562$ to $0.996761$.
\end{proposition}
The construction takes $S\sim\operatorname{Bernoulli}(0.2)$, with $Y\mid S=0$ uniform on $\{1,2\}$ and $Y\mid S=1$ uniform on $\{0,2\}$. At thresholds $(0,1)$, use $q=(0.1,0.5)$, $m(S)=\{0.5S,0.5+0.4(S-0.2)\}$ and independent verification with $p=0.1$. Figure~\ref{fig:separation} displays the two criteria. The proof uses marginal Gaussian probability inequalities, not the accuracy of a numerical optimiser.

\begin{figure}[htbp]\centering
\includegraphics[width=.82\linewidth]{figure1_separation.pdf}
\caption{The exact two-threshold example separates the trace and simultaneous-band criteria. Both curves are relative to Anchor; at the trace optimum $25/41$, the simultaneous critical value exceeds the anchor value. These are Gaussian calculations, not Monte Carlo estimates.}\label{fig:separation}
\noindent Alt text: The trace improves while the simultaneous criterion worsens at the trace optimum.
\end{figure}


\subsection{An implementable paired calibration}
We now replace the covariance set by a sample estimate. A second, independent sample $\sg$ estimates $\Omega$. We write $\wh F=\PP_{\sg}K_hU_0/\PP_{\sg}K_h$, $\wh r_i=(U_{0,i}-\wh F,B_i)^\top$ and
\begin{equation}\label{eq:moments}
 \wh\Omega=\frac p{v_h}\PP_{\sg}(K_h^2\wh r\wh r^\top),\quad
 \wh A=\frac p{v_h}\PP_{\sg}(K_h^2\wh r),\quad
 \wh\ell_i=\frac{K_{h,i}(U_{0,i}-\wh F)}{\PP_{\sg}K_h}.
\end{equation}
Equation~\eqref{eq:moments} gives the plug-in moments on the gate sample. The covariance influence contribution is
\begin{equation}\label{eq:covif}
 \wh C_i=\wh\Psi_i\wh\Psi_i^\top-\wh\Omega
       -E_0\wh\ell_i\wh A^\top-\wh A\wh\ell_i^\top E_0^\top,
       \qquad E_0=\binom{I_J}{0}.
\end{equation}
The last two terms in \eqref{eq:covif} account for estimating the anchor centring; omitting them changes the covariance uncertainty being calibrated.

Let $\Lset$ be a finite weight grid containing zero. For $\lambda>0$, define
\begin{equation}\label{eq:gradient}
 D_\lambda(\Omega)=\nabla_\Omega\{g_\Omega(\lambda)/\lambda\},\qquad
 \wh s_{i\lambda}=\inner{D_\lambda(\wh\Omega)}{\wh C_i}.
\end{equation}
We centre these scores across observations. Independent standard normal multipliers $\xi_i$ approximate
$\max\{0,\max_{\lambda>0}n_{\sg}^{-1}\sum_i\xi_i\wh s_{i\lambda}\}$.
Let $\wh t_\delta$ be an upper simulated quantile at probability $1-\delta$. We add a nonnegative second-order cushion $b_N$ and use fresh Gaussian draws to obtain simultaneous numerical intervals $[\underline c_\lambda,\overline c_\lambda]$ for the plug-in critical values. Note that the lower interval at the anchor is exactly the anchor critical value. We set
\begin{equation}\label{eq:rule}
 \wh L(0)=0,\quad
 \wh L(\lambda)=\underline c_0-\overline c_\lambda
                       -\frac{\wh t_\delta+b_N}{\sqrt{n_{\sg}}}\,\lambda
                       \quad(\lambda>0),\qquad
 \wh\lambda\in\argmax_{\lambda\in\Lset}\wh L(\lambda).
\end{equation}
Equation~\eqref{eq:rule} is the selection rule used throughout: a weight is adopted only when its certified lower bound on the gain is positive.

\begin{procedurebox}
\textbf{Algorithm 1: Certified local simultaneous inference}
\begin{enumerate}
\item Fix $x_0$, $h$, the threshold grid $\Yset_J$, the weight grid $\Lset$, the three roles and the error budgets $(0.049,0.0005,0.0005)$.
\item Fit the anchor $q$ and the candidate $m$ on the fitting role.
\item On the gate role, compute $\wh\Omega$, $\wh A$, the centred scores $\wh s_{i\lambda}$ and the multiplier quantile $\wh t_\delta$; form the numerical intervals $[\underline c_\lambda,\overline c_\lambda]$.
\item Select $\wh\lambda$ by \eqref{eq:rule}.
\item On the evaluation role, form $\wh F_{\se,\wh\lambda}$ and the simultaneous band with the calibrated critical value.
\item Rotate the three roles and aggregate the evaluation estimates.
\end{enumerate}
\end{procedurebox}


